%%%%%%%%%%%%%%%%%%%%%%%%%%%%%%%%%%%%
% Lesson Plan (50 minutes)
%%%%%%%%%%%%%%%%%%%%%%%%%%%%%%%%%%%%
\begin{frame}
    \frametitle{Lesson Plan}
    \begin{itemize}
        \item 4 min Lecture: Bayesian Knowledge Building --- restaurant analogy, prior/posterior updating
        \item 3 min Lecture: Bayes' Rule, Reinterpreted --- recall Bayes' rule for events, recall Module 5's parameter reinterpretation from Lecture 31
        \item 9 min Lecture: Disease Screening, Fully Worked --- full posterior derivation, why the posterior isn't higher, frequentist-vs-Bayesian comparison table
        \item 6 min Lecture: Sequential Updating --- second positive test, third test, ``yesterday's posterior becomes today's prior''
        \item 9 min Think/Pair/Share: \textit{Spam or Not Spam?} worksheet --- compute a posterior from one clue, then update again with a second, independent clue (see \texttt{Lecture33\_worksheet.tex})
        \item 12 min Lecture: More Than Two Hypotheses --- gravitational-wave classification, sequential updating across 4 hypotheses, recap table
        \item 2 min Lecture: Summary
    \end{itemize}
\end{frame}

%%%%%%%%%%%%%%%%%%%%%%%%%%%%%%%%%%%%
% Total: 45 min content + ~5 min buffer for agenda/LO/announcements and pacing slack = 50 min
% Activity placed right after Sequential Updating (2-hypothesis case) so students practice that
% skill themselves before the lecture escalates to the 4-hypothesis gravitational-wave example.
%%%%%%%%%%%%%%%%%%%%%%%%%%%%%%%%%%%%
